\section{Conclusion}\label{sec:conclusion65}
Neural Bellman Operators connect an attained own-future continuation to feasible decisions under the original economic law. The constructive neural realization, complete-policy comparison and explicit acquisition budget remain the paper's foundation. Under recursive risk, state-dependent qualification retains the original objective and the all-original-action lower obligation. The investment technology supplies primitive sufficient conditions under which these inequalities can actually be computed.

Convex-order risk integration adds a second-order precision account without smoothness of the risk-averse value. A complete admission-and-fallback service then turns the original-optimum guarantee into a support bound for paired economic inference. Together these results permit numerical resolution, sample size and a model-relative economic margin to be fixed before evaluation. This conclusion concerns the declared initial-state distribution and the original continuous innovation law, not an unreported change in either law.

The empirical questions remain distinct. A storage benefit need not be neural-specific; tighter uniform residuals need not yield better policy choices; and a valid screen need not save work after its construction is paid for. The new controlled-rank comparison and complete-service records make those distinctions observable. The acquisition-and-reuse region states which additional cost and price information supports an economic purchase. The fixed-accuracy horizon construction retains its explicit state, action, innovation and bit-growth costs; an unexecuted largest-vector comparison is not inferred from it.

The result is a constructive economic method with an original-law policy account and an explicit precision requirement. Its broader NBO applications, prior comparisons and adverse observations remain part of the preserved research record, under their own assumptions rather than an unsupported universal superiority assertion.
